\section{问题重述}
题目将算子细化为由原子指令构成的有向无环图，节点包括指定执行单元上的操作，以及缓冲区申请和释放操作\cite{statement}。每个缓冲区只属于一种缓存，需占据长度确定的连续地址区间。不同执行单元可以并行工作，同一执行单元只能串行执行。数据依赖和物理地址复用都会引入等待，因此拓扑顺序、缓存位置与实际完成时间必须联合考虑。

问题一要求在不考虑具体缓存容量限制时，寻找最大累计驻留量尽可能小的拓扑序。问题二在给定各类缓存容量后，为缓冲区确定偏移，必要时加入换出及换入操作，使额外搬运量尽可能小。问题三则允许进一步调整顺序和分配，在搬运量不显著增加的前提下降低执行时间。三个问题分别突出存储压力、物理可行性和执行并行性，目标之间存在竞争。

本文对六个输入图使用相同算法组合，不以算子名称编写专用调度规则。题面未给出“不显著增加”的数值阈值，因此问题三采用更严格的搬运量零增加约束。研究重点是获得可核验的完整调度方案，并明确启发式与评价模型的适用边界。
